\title[Intrinsic reconstruction of webs]{Intrinsic reconstruction of webs from nonreduced multiplication-failure schemes}
\author{Qian Qi}
\date{September 23, 2026 \quad (Revision 136)}
\subjclass[2020]{14M12, 14J28, 13C40, 14B05, 14D20, 05E05}
\keywords{Intrinsic reconstruction, webs of quadrics, Fitting schemes, Schur functors, exterior support, Jacobian quartics, nilpotent filtration}
\hypersetup{pdftitle={Intrinsic reconstruction of webs from nonreduced multiplication-failure schemes, revision 136},pdfauthor={Qian Qi}}
\begin{abstract}
We prove that the abstract nonreduced multiplication-failure scheme
of a cube-zero algebra with Hilbert function \((1,4,6)\) determines
its relation web up to projective equivalence whenever the web is
basepoint-free and its Jacobian quartic is smooth.  No further
genericity condition is required.  The intrinsic deepest normal cone
recovers both Pluecker components, transported by one projective
transformation.  Its residual coefficient module is recovered by a
graded ideal quotient, with the determinant line explicitly retained.
The support calculation extends to every dimension: we determine the
kernel of contraction on the determinant-twisted degree-\(n\)
Jacobian component in \(\bigwedge^n\Sym^2V\) for every bilinear rank.
An algebraic cofactor identity gives the exceptional even-dimensional
kernel.  For every \(n\ge4\), three essential variables force
exterior support greater than \(2n\), yielding reconstruction from
the two complementary component lines on a nonempty open family,
compatibly with arbitrary base change.  In dimension four this
excludes every exceptional recombination in the smooth geometric
family.  The first nilpotent layer recovers its polarized Jacobian
K3, while the full scheme separates the remaining web classes.
A separate supplement retains the primary-boundary and specialization
theory.
\end{abstract}
\maketitle
